\subsection{Theorem-Style Results and Formal Statements}

Formal environments identify a block's function, evidential status, and permissible use as authority.

\begin{convention} \label{conv:theorem-style-result-hierarchy} \hlt{Theorem-Style Result Hierarchy}\\
In these notes, the theorem-style result environments are exactly the following; each contains a proved claim.
\begin{enumerate}[label=(\roman*)]
\setlength{\itemsep}{0pt}
\item \hlt{Theorem}: a major result that advances a principal objective of the surrounding work.
\item \hlt{Proposition}: an independently useful result of smaller local weight.
\item \hlt{Lemma}: a subsidiary result introduced primarily to support a later proof.
\item \hlt{Corollary}: a direct consequence of an earlier proved result that requires little additional argument.
\end{enumerate}
These names record local structural roles, not a universal ranking of importance.\\
Definitions stipulate objects or terminology, whereas conjectures record unproved claims.\\
Convention follows standard mathematical-writing guidance \citep{mit1998mathematicalwriting}.
\end{convention}

\begin{process} \label{proc:theorem-statement-audit} \hlt{Theorem-Statement Construction and Audit}\\
Given a candidate conclusion and its supporting argument, produce a self-contained theorem-style statement that can be cited and proved without hidden context.
\begin{enumerate}[label=(\roman*)]
\setlength{\itemsep}{0pt}
\item \hlt{Purpose}: isolate the reusable conclusion and identify its role in the dependency chain.
\item \hlt{Objects}: bind every ambient set, structure, parameter, map, and index before its first use.
\item \hlt{Hypotheses}: retain only the argument's exact domain, regularity, finiteness, independence, and boundary conditions.
\item \hlt{Conclusion}: state existence, uniqueness, characterization, bound, convergence, or other claim explicitly.
\item \hlt{Claim census}: itemize each conclusion whose omission would remove a distinct proof obligation.
\item \hlt{Role}: select theorem, proposition, lemma, or corollary by Convention \ref{conv:theorem-style-result-hierarchy}.
\item \hlt{Locality}: move motivation, examples, and proof machinery outside the statement, and delete unused assumptions.
\item \hlt{Traceability}: attach a stable title and label, and cite the external owner of a nonoriginal result.
\end{enumerate}
The audit passes only when the statement remains intelligible apart from the preceding paragraph and every symbol is defined at first use. These are central expository checks \citep{knuth1988mathematicalwriting}.
\end{process}

\begin{principle} \label{princ:claim-authority-separation} \hlt{Claim--Authority Separation}\\
Match each formal block's environment and evidence to its epistemic status. A definition stipulates meaning; an assumption declares a condition; a conjecture remains unproved; a theorem-style result requires a proof; and an empirical statement requires data and a protocol. Only proved theorem-style results may serve as mathematical proof authority. A citation records prior ownership but does not discharge a proof obligation when the present work claims to establish the result.
\end{principle}

\begin{warning} \label{warn:statement-proof-contamination} \hlt{Statement--Proof Contamination}\\
If a theorem-style statement contains a proof plan, auxiliary construction, example, or claim of obviousness, its logical contract becomes obscure and later citation becomes unreliable. Move proof-specific objects into the proof, place explanatory material after the formal owner, and delete hypotheses unused by the proof.
\end{warning}
